\chapter{Conditionals, Iterations, Functions, and Recursion}
\label{ch:control_functions}

\section{Control Flow and Conditional Branching}

Programs make decisions by evaluating boolean expressions to alter the sequence of instruction execution. Python uses strict indentation blocks (conventionally 4 spaces) rather than curly braces to define scope boundaries.

\subsection{Syntax Structures: \texttt{if}, \texttt{if-else}, and \texttt{if-elif-else}}
\begin{lstlisting}[caption={Multi-way Signal Classification}]
signal_level_dbm = -45.0

if signal_level_dbm >= -30.0:
    quality = "Strong Signal (Optimal)"
elif signal_level_dbm >= -60.0:
    quality = "Good Signal (Acceptable)"
elif signal_level_dbm >= -80.0:
    quality = "Weak Signal (Marginal)"
else:
    quality = "No Signal / Link Failure"

print(f"Receiver Status: {quality}")
\end{lstlisting}

\subsection{Short-Circuit Evaluation}
Python evaluates logical operators (\texttt{and}, \texttt{or}) using short-circuit semantics:
\begin{itemize}[leftmargin=1.8em, itemsep=2pt]
    \item \textbf{\texttt{A and B}:} If $A$ is false, Python returns $A$ immediately without evaluating $B$.
    \item \textbf{\texttt{A or B}:} If $A$ is true, Python returns $A$ immediately without evaluating $B$.
\end{itemize}
This allows safe guard expressions, such as \texttt{if denom != 0 and (num / denom) > threshold:}.

\section{Iterative Control Structures}

\subsection{The \texttt{while} Loop and Sentinel Controls}
The \texttt{while} statement executes a suite of instructions repeatedly as long as a test condition remains true. It is well-suited for iterative numerical approximations (e.g., Newton-Raphson method) where the iteration count is not known a priori.

\begin{lstlisting}[caption={Newton-Raphson Square Root Approximation}]
def newton_sqrt(n, tolerance=1e-7):
    if n < 0:
        raise ValueError("Cannot compute real square root of negative number.")
    if n == 0:
        return 0.0
    x = n / 2.0
    while True:
        root_next = 0.5 * (x + n / x)
        if abs(root_next - x) < tolerance:
            return root_next
        x = root_next

print("Calculated:", newton_sqrt(2.0))
\end{lstlisting}

\subsection{The \texttt{for} Loop and the \texttt{range()} Sequence Generator}
Python \texttt{for} loops iterate over elements of any sequence or iterable protocol:
\begin{equation}
\texttt{range(start, stop, step)} \implies [start, start+step, start+2\cdot step, \dots < stop]
\end{equation}
Python supports loop control statements: \texttt{break} (terminates loop), \texttt{continue} (skips current iteration), and an optional \texttt{else} block that executes if and only if the loop terminates normally without encountering a \texttt{break}.

\section{Modular Architecture: Functions}

Functions encapsulate logic, promote reusability, and establish isolated local scopes.

\subsection{Parameter Binding Semantics}
Python functions support diverse parameter passing mechanisms:
\begin{enumerate}[leftmargin=1.8em, itemsep=3pt]
    \item \textbf{Positional Arguments:} Matched strictly left-to-right.
    \item \textbf{Keyword Arguments:} Specified by name, independent of argument order (\texttt{func(b=2, a=1)}).
    \item \textbf{Default Parameters:} Fallback values evaluated once at definition time.
    \item \textbf{Arbitrary Positional Arguments (\texttt{*args}):} Bundles excess positional arguments into a \texttt{tuple}.
    \item \textbf{Arbitrary Keyword Arguments (\texttt{**kwargs}):} Bundles excess keyword arguments into a \texttt{dict}.
\end{enumerate}

\begin{warningbox}[Mutable Default Arguments Pitfall]
Default arguments are evaluated once when the function definition is executed, \textbf{not} at invocation time. Never use mutable objects (like lists or dictionaries) as default arguments:
\begin{lstlisting}
# BUG: The list accumulator persists across invocations!
def append_reading(val, readings=[]):
    readings.append(val)
    return readings

# CORRECT IDIOM:
def append_reading_safe(val, readings=None):
    if readings is None:
        readings = []
    readings.append(val)
    return readings
\end{lstlisting}
\end{warningbox}

\subsection{Scope Resolution: The LEGB Rule}
Identifier lookup follows four nested namespaces:
\begin{enumerate}[leftmargin=1.8em, itemsep=2pt]
    \item \textbf{L (Local):} Names assigned inside current function.
    \item \textbf{E (Enclosing):} Names in the enclosing functions (closures).
    \item \textbf{G (Global):} Names assigned at the top-level of module file.
    \item \textbf{B (Built-in):} Names pre-assigned in Python's \texttt{builtins} module.
\end{enumerate}
Modifying variables outside local scope requires the \texttt{global} or \texttt{nonlocal} keywords.

\section{Functional Paradigms and Anonymous Functions}
Python provides \texttt{lambda} expressions for lightweight, anonymous single-expression functions:
\begin{lstlisting}[caption={Higher-Order Functional Tools}]
from functools import reduce

voltages = [1.2, 3.3, 5.0, 12.0, 24.0]

# map: transform each element
scaled = list(map(lambda v: v * 1.05, voltages))

# filter: extract subset matching predicate
high_voltages = list(filter(lambda v: v >= 5.0, voltages))

# reduce: combine elements recursively
total_energy = reduce(lambda acc, v: acc + v**2, voltages, 0)
\end{lstlisting}

\section{Recursion and Memory Stack Mechanics}

A recursive function calls itself to solve smaller subproblems until reaching a base condition.

\begin{definition}[Call Stack Frame Allocation]
Every recursive call pushes an activation record (stack frame) containing local variables, parameters, and return program counters onto the execution call stack. If recursion depth exceeds \texttt{sys.getrecursionlimit()} (typically 1000), a \texttt{RecursionError} is raised to prevent stack overflow crashes.
\end{definition}

\begin{examplebox}[Binary Search Implementation via Recursion]
\begin{lstlisting}
def binary_search(arr, target, low, high):
    if low > high:
        return -1  # Base case: element not found
    
    mid = (low + high) // 2
    if arr[mid] == target:
        return mid
    elif arr[mid] > target:
        return binary_search(arr, target, low, mid - 1)
    else:
        return binary_search(arr, target, mid + 1, high)
\end{lstlisting}
Time complexity: $\mathcal{O}(\log n)$, Space complexity: $\mathcal{O}(\log n)$ due to recursive stack frames.
\end{examplebox}

\section{Chapter Summary and Review Questions}
\begin{enumerate}[leftmargin=1.8em, itemsep=4pt]
    \item Contrast \texttt{break}, \texttt{continue}, and \texttt{pass} with practical code examples.
    \item How does Python resolve variable identifiers under the LEGB hierarchy? Give an example where \texttt{nonlocal} is required.
    \item Write a recursive function to compute the Greatest Common Divisor (GCD) of two integers using the Euclidean algorithm.
    \item Explain the consequence of defining \texttt{def register(item, log=[])} and suggest the standard defensive fix.
\end{enumerate}
